% 同一边集附加同一组数值：{1,2}:2, {1,3}:1, {2,3}:3, {2,4}:1, {3,4}:4。
% 长度读法：路径1-2-4成本2+1=3，路径1-3-4成本1+4=5。下排对照无边与零权边。
\begin{tikzpicture}[>=Stealth,
  v/.style={circle,draw=FigureStroke,fill=FigureBlueFill,line width=.8pt,minimum size=6.5mm,inner sep=0pt},
  e/.style={draw=FigureStroke,line width=.8pt},
  w/.style={fill=white,inner sep=1pt,text=FigureInk},
  note/.style={font=\FigureNoteFont\fontsize{7.5}{9}\selectfont,text=FigureMuted,align=center}]
\begin{scope}[xshift=0cm]
  \node[v] (1) at (-1.2,0) {$1$}; \node[v] (2) at (0,.85) {$2$};
  \node[v] (3) at (0,-.85) {$3$}; \node[v] (4) at (1.2,0) {$4$};
  \draw[e] (1)--node[w]{$2$}(2);
  \draw[e] (1)--node[w]{$1$}(3);
  \draw[e] (2)--node[w]{$3$}(3);
  \draw[e] (2)--node[w]{$1$}(4);
  \draw[e] (3)--node[w]{$4$}(4);
  \node at (0,1.45) {(a) 长度$\ell$};
  \node[note] at (0,-1.5) {沿路径相加：\\$1,2,4$的成本为$3$};
\end{scope}
\begin{scope}[xshift=3.7cm]
  \node[v] (1) at (-1.2,0) {$1$}; \node[v] (2) at (0,.85) {$2$};
  \node[v] (3) at (0,-.85) {$3$}; \node[v] (4) at (1.2,0) {$4$};
  \draw[e] (1)--node[w]{$2$}(2);
  \draw[e] (1)--node[w]{$1$}(3);
  \draw[e] (2)--node[w]{$3$}(3);
  \draw[e] (2)--node[w]{$1$}(4);
  \draw[e] (3)--node[w]{$4$}(4);
  \node at (0,1.45) {(b) 容量$c$};
  \node[note] at (0,-1.5) {限制通过量：\\$\{1,3\}$上至多通过$1$};
\end{scope}
\begin{scope}[xshift=7.4cm]
  \node[v] (1) at (-1.2,0) {$1$}; \node[v] (2) at (0,.85) {$2$};
  \node[v] (3) at (0,-.85) {$3$}; \node[v] (4) at (1.2,0) {$4$};
  \draw[e] (1)--node[w]{$2$}(2);
  \draw[e] (1)--node[w]{$1$}(3);
  \draw[e] (2)--node[w]{$3$}(3);
  \draw[e] (2)--node[w]{$1$}(4);
  \draw[e] (3)--node[w]{$4$}(4);
  \node at (0,1.45) {(c) 相似度$w$};
  \node[note] at (0,-1.5) {大权重更强地\\要求两端取值接近};
\end{scope}
% 下排：无边与零权边
\begin{scope}[yshift=-3.2cm,xshift=.5cm]
  \node[v] (a) at (0,0) {$1$}; \node[v] (b) at (1.4,0) {$2$};
  \node at (.7,-.7) {(d) 无边};
\end{scope}
\begin{scope}[yshift=-3.2cm,xshift=3.6cm]
  \node[v] (a) at (0,0) {$1$}; \node[v] (b) at (1.4,0) {$2$};
  \draw[e] (a)--node[w]{$0$}(b);
  \node at (.7,-.7) {(e) 零权边};
\end{scope}
\node[note,anchor=west] at (6.1,-3.2) {两者的数值条目都是$0$，\\只有(e)中$1,2$相邻};
\end{tikzpicture}
